\par\noindent\SourceLeaf{55}%
\hspace*{\parindent}Sic per coefficientes binomiales brevissime exprimi posse videtur lex complicatissima.
Licet enim, evolutis, ut fieri potest, in producta his coefficientibus, factores alterius
nonnulli tollantur ab alterius factoribus, tamen expressio finitima longe prolixior fit,
ac prior fuerat. Sic enim, ut facto rei veritas comprobetur, brevissime exhibetur illa
ratione coefficientis indefiniti $r^{\text{ti}}$ terminus generalis $h^{\text{tus}}$

\[
\frac{\nu.(\nu-1).\,.\,.(\nu-(r+2h-1))}{1.2^2.3^2.\,.\,. h^2} . \frac{1}{(h+1)(h+2).\,.\,.(h+r)}\frac{1}{2^{r+2h-1}}n^{r+2h}.
\]

\textls[120]{Adnotatio.} 
Eiusmodi formulam sortietur, qui ex seriebus pro coefficientibus aliquot initialibus ab
\textsc{Eulero}\HistoricalNote{*)}{*)}{\textsc{Euler} l. c. Probl. 36.} propositis legem generalem abstrahere velit. Licet via, quam ad scrutandam
legem istam indicat Auctor Illustris, ad obtinendas series iamiam expositas minus
directa ac facilis videatur, alio tamen respectu praeclara, et minime negligenda videtur.
Omne artificium in eo consistit, ut ficta pro quantitate $(1+n\cos\varphi)^{\nu}$ serie

\[
A+A^{(1)}\cos\varphi+A^{(2)}\cos2\varphi+.\,.\,.\,.
\]

\noindent per differentiationem eruatur scala recursionis, qua istius coefficientes a se invicem
pendent. Quae, cum tantummodo tripartita sit, per binos quosvis antecedentes
coefficientem quemvis exhibet, ut igitur, modo priores bini finiti fuerint, ceteri omnes
forma finita exhiberi queant. Quae mutua coefficientium relatio, licet ex seriebus
iamiam inventis erui possit, facilius tamen artificio \textsc{Euler}iano evincitur.

\Needspace{8\baselineskip}
7. Sit, ut antea

\[
(1+n\cos\varphi)^{\nu}=A+A^{(1)}\cos\varphi+A^{(2)}\cos2\varphi+.\,.\,.\,.+A^{(h)}\cos h\varphi+.\,.\,.
\]

\noindent Sumtis ab utraque parte logarithmorum differentialibus, habebitur, remoto divisione
ipso $d\varphi$

\[
\frac{\nu n\sin\varphi}{1+n\cos\varphi}=\frac{A^{(1)}\sin\varphi+2A^{(2)}\sin2\varphi+3A^{(3)}\sin3\varphi+.\,.\,.+hA^{(h)}\sin h\varphi+.\,.\,.}{A+A^{(1)}\cos\varphi+A^{(2)}\cos2\varphi+.\,.\,.+A^{(h)}\cos h\varphi+.\,.\,.}
\]

$\alpha$) Sublatis multiplicando utrimque divisoribus habetur ex una parte

\[
\begin{aligned}
&\nu n\sin\varphi\bigl(A+A^{(1)}\cos\varphi+A^{(2)}\cos2\varphi+.\,.\,.\,.+A^{(h)}\cos h\varphi\\&\quad+A^{(h+1)}\cos(h+1)\varphi+A^{(h+2)}\cos(h+2)\varphi+.\,.\,.\,.\bigr).
\end{aligned}
\]
